% Successor v10: first full-prefix difference, 10 October 2026.
% Review-only TODO fragment. Replace the older Lemma 6.49 TODO with this
% information when the four mathematical patches are assembled. This
% is not a patch verified against the unavailable v10 main.tex archive.
%
% Candidate addition at Lemma 6.49:
\todo[inline]{\textbf{v10 verification (converse record step).}
For two original types from the same enumerated partial structure,
whose first $d$-level $L^+$ prefixes agree and whose types exist
through level $d+1$, a discrepancy of their complete prefixes at
$d+1$ is equivalent to a discrepancy in one of the two directed
binary relation vectors connecting their type vertices to the
new socle coordinate $d$.
The predecessor agreement itself implies the common singleton root;
one need not add the root equality as an independent hypothesis.
All $E$-pairs are determined because $d+1$ is bounded by both
free levels. This equivalence is formalised in
\texttt{V10/FirstFullTypeDifference.lean}, using
\texttt{PartialTypeRestriction.lean} and
\texttt{CommonSocleType.lean}.
To finish the meet lemma, still identify the actual $K^\mathcal F_{pt}$
ancestor of a node at level $d$ with the complete extracted prefix,
and show that the numerical $H$ traces faithfully encode \emph{both}
binary orientations, including odd fake socle vertices.
In particular, do not replace the sharp non-top condition
$p+1<j$ by $p<j$: the candidate positive mismatch
$\iota(p,m)$ must lie strictly below
$\mathrm{fl}_{H^+}(\iota(j,m))$ to occur inside its partial type.}
%
% Endpoints, mark as proved only after exact SHA's green CI:
%   sameAmbient_root_of_prefix_eq
%   sameAmbient_first_type_difference_has_binary_witness
%   sameAmbient_first_type_difference_iff_binary
